\documentclass[11pt]{article}

% ===================== PACKAGES =====================
\usepackage[a4paper,margin=1in]{geometry}
\usepackage{amsmath,amssymb}
\usepackage{enumitem}
\usepackage{fancyhdr}
\usepackage{xcolor}

% ===================== HEADER & FOOTER =====================
\pagestyle{fancy}
\fancyhf{}
\lhead{CSE400: Fundamentals of Probability in Computing}
\rhead{Lecture 20}
\cfoot{\thepage}

% ===================== DOCUMENT =====================
\begin{document}

\begin{center}
    {\Large \textbf{Lecture 20: Linear Transformations and Expectations of Multiple Random Variables}}
\end{center}

\vspace{0.5cm}

% ===================== SECTION 1 =====================
\section{Overview}

This lecture covers:
\begin{itemize}
    \item Statistical representation of random vectors
    \item Mean vector, correlation, and covariance matrices
    \item Linear transformations
    \item Mapping properties under $Y = AX + b$
    \item Multivariate Gaussian distributions
    \item Extension from 1D/2D to N-dimensional joint PDFs
    \item Special cases and properties
\end{itemize}

% ===================== SECTION 2 =====================
\section{Statistical Representation of Random Vectors}

A system involving multiple random variables is represented as a random vector.

Let:
\[
X = [X_1, X_2, \dots, X_N]^T
\]

% ===================== SECTION 3 =====================
\section{Case Study: Smart City Traffic + Weather System}

An intelligent system is designed to predict traffic congestion using transportation and climate variables.

\subsection*{Step 1: Define N Random Variables}

A vector of random variables is defined:
\begin{itemize}
    \item Represents system variables (traffic, weather, etc.)
    \item Forms an N-dimensional random vector ($N = 5$)
\end{itemize}

\subsection*{Step 2: Mean Vector}

The mean vector represents typical or expected values of each variable.

Let:
\[
\mu = E[X]
\]

\subsection*{Step 3: Covariance Matrix}

The covariance matrix captures dependencies between variables.

Interpretation:
\begin{itemize}
    \item More rain $\Rightarrow$ more congestion (traffic increases)
    \item More rain $\Rightarrow$ lower speed
    \item More traffic $\Rightarrow$ higher pollution
    \item Higher temperature $\Rightarrow$ pollution disperses
\end{itemize}

% ===================== SECTION 4 =====================
\section{Joint Gaussian Distribution}

\textbf{Assumption:} The system follows a joint Gaussian distribution.

This implies:
\begin{itemize}
    \item The system is fully described by:
    \begin{itemize}
        \item Mean vector
        \item Covariance matrix
    \end{itemize}
\end{itemize}

% ===================== SECTION 5 =====================
\section{Linear Transformations}

A new variable is defined using a linear transformation:
\[
Y = AX + b
\]

\subsection*{Congestion Index}

The city defines a Congestion Index as a linear transformation of the variables.

% ===================== SECTION 6 =====================
\section{Inference from Transformation}

\textbf{Inference \#1}

The transformation maps the original variables into a new representation (Congestion Index).

\textbf{Inference \#2}

If rainfall increases suddenly:
\begin{itemize}
    \item Speed $\downarrow$
    \item Traffic $\uparrow$
    \item Pollution $\uparrow$
\end{itemize}

% ===================== SECTION 7 =====================
\section{Key Takeaways}

\begin{itemize}
    \item Random vectors model multiple random variables jointly
    \item Mean vector captures expected values
    \item Covariance matrix captures relationships between variables
    \item Joint Gaussian distribution fully described by mean and covariance
    \item Linear transformations map random vectors into new variables
    \item System behavior can be inferred through dependencies and transformations
\end{itemize}

\end{document}